\pdfoutput=1
\documentclass[11pt,a4paper,reqno]{amsart}
\usepackage{amssymb}

\usepackage{bm}
\numberwithin{equation}{section}

\usepackage{mathabx}
\usepackage{mathtools}%                  http://www.ctan.org/pkg/mathtools
%\usepackage{color}



\theoremstyle{plain}

\newtheorem{theorem}[equation]{Theorem}
\newtheorem{proposition}[equation]{Proposition}
\newtheorem{lemma}[equation]{Lemma}
\newtheorem{corollary}[equation]{Corollary}
\newtheorem{conjecture}[equation]{Conjecture}
\newtheorem{principle}[equation]{Principle}
\newtheorem{claim}[equation]{Claim}

\theoremstyle{definition}

\newtheorem{definition}[equation]{Definition}
\newtheorem{remark}[equation]{Remark}
\newtheorem{remarks}[equation]{Remarks}
\newtheorem{example}[equation]{Example}
\newtheorem{examples}[equation]{Examples}
%\newtheorem{problem}[subsection]{Problem}
%\newtheorem{question}[subsection]{Question}

\newcommand\R{\mathbb{R}}
\newcommand\Z{\mathbb{Z}}
\newcommand\N{\mathbb{N}}
\newcommand\Q{\mathbb{Q}}
\newcommand\eps{\varepsilon}
\newcommand\normsymb{\ell}  % was \|\,\|
\newcommand\norm[1]{\ell({#1})}   % was \|{#1}\|
\newcommand\normn[1]{\ell_n({#1})}   % was \|{#1}\|
\newcommand\free[1]{\mathbf{F}_{#1}}


\usepackage[breaklinks=true,hidelinks]{hyperref}


\begin{document}

\title{Homogeneous length functions on groups}

\author{D.H.J. Polymath}
\address{\url{http://michaelnielsen.org/polymath1/index.php}}
%\email{}

\begin{abstract}
A pseudo-length function defined on an arbitrary group $G = (G,\cdot,e,
(\,)^{-1})$ is a map $\normsymb: G \to [0,+\infty)$ obeying $\norm{e}=0$,
the symmetry property $\norm{x^{-1}} = \norm{x}$, and the triangle
inequality $\norm{xy} \leq \norm{x} + \norm{y}$ for all $x,y \in G$. We
consider pseudo-length functions which saturate the triangle inequality
whenever $x=y$, or equivalently those that are \emph{homogeneous} in the
sense that $\norm{x^n} = n\,\norm{x}$ for all $n\in\N$. We show that this
implies that $\norm{[x,y]}=0$ for all $x,y \in G$. This leads to a
classification of such pseudo-length functions as pullbacks from
embeddings into a Banach space.  We also obtain a quantitative version of
our main result which allows for defects in the triangle inequality or
the homogeneity property.
\end{abstract}

\subjclass[2010]{20F12, 20F65}

\keywords{homogeneous length function,
pseudo-length function,
quasi-morphism,
Banach space embedding}

\maketitle
%\today


%{{{1 Section 1 - Introduction
\section{Introduction}

Let $G = (G, \cdot, e, (\,)^{-1})$ be a group (written multiplicatively,
with identity element $e$).  A \emph{pseudo-length function} on $G$ is a
map $\normsymb : G \to [0,+\infty)$ that obeys the properties
\begin{itemize}
\item $\norm{e} = 0$,
\item $\norm{ x^{-1} } = \norm{x}$,
\item $\norm{ x y } \leq \norm{x} + \norm{ y }$
\end{itemize}
for all $x,y\in G$.  If in addition we have $\norm{x}>0$ for all $x \in G
\setminus \{e\}$, we say that $\normsymb$ is a \emph{length function}. By
setting $d(x,y) \coloneqq \norm{x^{-1}y}$, it is easy to see that
pseudo-length functions (resp. length functions) are in bijection with
left-invariant pseudometrics (resp. left-invariant metrics) on $G$.

From the axioms of a pseudo-length function it is clear that one has the
upper bound
\[
\norm{x^n} \leq |n| \, \norm{x}
\]
for all $x \in G$ and $n \in \Z$.  Let us say that a pseudo-length
function $\normsymb: G\to [0,+\infty)$ is \emph{homogeneous} if equality
is always attained here, thus one has
\begin{equation}\label{linear-growth}
\norm{x^n} = |n| \, \norm{x}
\end{equation}
for all $x \in G$ and any $n \in \Z$.  Using the axioms of a
pseudo-length function, it is not difficult to show that the homogeneity
condition \eqref{linear-growth} is equivalent to the triangle inequality
holding with equality whenever $x=y$ (i.e., that \eqref{linear-growth}
holds for $n=2$); see \cite[Lemma 1]{GK}.

If one has a real or
complex Banach space $\mathbb{B} = (\mathbb{B},\| \, \|)$, and $\phi: G
\to \mathbb{B}$ is any homomorphism from $G$ to $\mathbb{B}$ (viewing the
latter as a group in additive notation), then the function $\normsymb: G
\to [0,+\infty)$ defined by $\norm{x} \coloneqq \| \phi(x) \|$ is easily
verified to be a homogeneous pseudo-length function.  Furthermore, if
$\phi$ is injective, then $\normsymb$ will in fact be a homogeneous
length function.  For instance, the function
\[
\norm{ (n, m) } \coloneqq |n + \sqrt{2} m|
\]
is a length function on $\Z^2$, where in this case $\mathbb{B} \coloneqq
\R$ and $\phi(n,m) \coloneqq n  + \sqrt{2} m$.  On the other hand, one
can easily locate many length functions that are not homogeneous, for
instance by taking the square root of the length function just
constructed.

The main result of this paper is that such Banach space constructions are
in fact the \emph{only} way to generate homogeneous (pseudo-)length
functions.

\begin{theorem}[Classification of homogeneous length functions]\label{class}
Given a group $G$, let $\normsymb: G \to [0,+\infty)$ be a homogeneous
pseudo-length function.  Then there exist a real Banach space $\mathbb{B}
= (\mathbb{B},\| \, \|)$ and a group homomorphism $\phi : G \to
\mathbb{B}$ such that $\norm{x} = \| \phi(x) \|$ for all $x \in G$.
Furthermore, if $\normsymb$ is a length function, one can take $\phi$ to
be injective, i.e., an isometric embedding.
\end{theorem}

We will derive Theorem~\ref{class} from a more quantitative result
bounding the pseudo-length of a commutator
\begin{equation}\label{comdef}
[x,y] \coloneqq xyx^{-1}y^{-1};
\end{equation}
see Proposition~\ref{main} below.  Our arguments will be elementary,
relying on directly applying the axioms of a homogeneous length function
to various carefully chosen words of $x$ and $y$, and repeatedly taking
an asymptotic limit $n \to \infty$ to dispose of various error terms that
arise in the estimates obtained in this fashion.

As one quick corollary of Theorem \ref{class}, we will obtain the
following characterization of the groups that admit
homogeneous length functions.

\begin{corollary}\label{abel}
A group admits a homogeneous length function if and only if it is
abelian and torsion-free.
\end{corollary}
%}}}


%{{{1 Section 1.1 - Examples and approaches
\subsection{Examples and approaches}

We now make a few remarks to indicate the nontriviality of
Theorem~\ref{class}. Corollary \ref{abel} asserts in particular that
there are no non-abelian groups with homogeneous length functions.
Whether or not such a striking geometric rigidity phenomenon holds was
previously unknown to experts. Moreover, the corollary fails
to hold if one or more of the precise conditions in the theorem are
weakened. For instance, such pseudo-length functions indeed exist (i)~on
non-abelian monoids, and (ii)~on all balls of finite radius in free
groups. We explain these two cases further in Section \ref{remarks-sec}.

Given these cases, one could \textit{a priori} ask if every non-abelian
group admits a homogeneous length function. This is not hard to
disprove; here are two examples.

\begin{example}[Nilpotent groups]
If $G$ is a nilpotent group of nilpotency class $2$ (such as the
Heisenberg group), then one has the identity $[x,y]^{n^2} = [x^n,y^n]$
for all $x,y \in G$ and natural numbers $n$ (since the map $(g,h) \mapsto
[g,h]$ is now a bihomomorphism $: G \times G \to [G,G] \subset Z(G)$).
If $[x,y]$ is non-trivial, then any homogeneous length function on $G$
would assign a linearly growing quantity to the right-hand side and a
quadratically growing quantity to the left-hand side, which is absurd;
thus such groups cannot admit homogeneous length functions.  The claim
then also follows for nilpotent groups of higher nilpotency class, since
they contain subgroups of nilpotency class $2$.
\end{example}


\begin{example}[Connected Lie groups]
As we explain in Remark~\ref{Rinv}, a homogeneous length function in fact
induces a \textit{bi-invariant} metric. Thus, if $G$ is a connected Lie
group with a homogeneous length function, then by \cite[Lemma
7.5]{Milnor}, $G \cong K \times \R^n$ for some compact Lie group $K$ and
integer $n \geq 0$. By \eqref{linear-growth}, $K$ cannot have torsion
elements, hence must be trivial. But then $G$ is abelian.
\end{example}

Prior to Corollary \ref{abel}, the above examples left open the question
of whether \textit{any} non-abelian group admits a homogeneous length
function. One may as well consider groups generated by two non-commuting
elements. As a prototypical example, let $\free{2}$ be the free group on
two generators $a,b$.  The word length function on $\free{2}$ is a length
function, but it is not homogeneous, since for instance the word length
of $(bab^{-1})^n = ba^n b^{-1}$ is $n+2$, which is not a linear function
of $n$.  (It is however the case that the word length of $x^n$ has linear
\emph{growth} in $n$ for any non-trivial $x$.) Similarly for the
Levenshtein distance (edit distance) on $\free{2}$.

Our initial attempts to construct homogeneous length functions on
$\free{2}$ all failed in light of our main result. However, many of these
methods apply under minor weakening of the hypotheses, such as working
with monoids rather than groups, or weakening homogeneity. Results in
these cases are discussed further in Section \ref{remarks-sec}.
%}}}


%{{{1 Section 1.2 - Further motivations
\subsection{Further motivations}

We conclude this part with some motivations from functional analysis and
probability, or more precisely the study of Banach space embeddings. If
$G$ is an additive subgroup of a Banach space $\mathbb{B}$, then clearly
the norm on $\mathbb{B}$ restricts to a homogeneous length function on
$G$. In~\cite{CSC,GK} one can find several equivalent conditions for a
given length function on a given group to arise in this way (studied in
the broader context of additive mappings and separation theorems in
functional analysis). These conditions are summarized in the following:

\begin{proposition}[{see \cite[Proposition 3.9]{KR2}}]\label{Pnormed}
Given a group $G$ equipped with a length function $\normsymb : G \to
[0,+\infty)$, the following are equivalent:
\begin{itemize}
\item[\textup{(1)}] There exists a Banach space $\mathbb{B}$ and a group
homomorphism $G \to \mathbb{B}$ that is an isometric embedding.

\item[\textup{(2)}] $G$ is abelian and $d_G(x,y) \coloneqq \norm{ x^{-1}
y }$ is a translation-invariant metric for which $G$ is `normed',
i.e., $d_G(e,x^n) = |n|\, d_G(e,x)$ for each $x\in G$ and $n\in \Z$. In
particular, $G$ is torsion-free.

\item[\textup{(3)}] $G$ is $\{ 2 \}$-normed \textup{$($}i.e.,
$\norm{ x^2 } = 2 \norm{ x }$ for each $x\in G$\textup{$)$} and is {\em
weakly commutative}, i.e., for all $x,y \in G$ there exists $n = n(x,y)
\in \mathbb{N}$ such that $(xy)^{2^n} = x^{2^n} y^{2^n}$.

\item[\textup{(4)}] $G$ is $\{ 2 \}$-normed and amenable.
\end{itemize}
\end{proposition}

In view of these equivalences, it is natural to try and characterize the
groups possessing a norm. This question is answered by
Corollary~\ref{abel}, which shows these are precisely the abelian
torsion-free groups. Thus in assertion (2), the hypothesis of admitting a
`norm' (i.e., homogeneous length function) is equivalent to that of being
abelian and torsion-free; and hence one may use either set of data to
infer the remaining assertions in Proposition \ref{Pnormed}.

Groups and semigroups with translation-invariant metrics also naturally
arise in probability theory, with the most important `normed' (i.e.,
homogeneous) examples being Banach spaces~\cite{LT}. Notice however that
in certain fundamental stochastic settings, formulating -- and proving --
results does not require the full Banach space structure. In this vein, a
general variant of the Hoffmann-J{\o}rgensen inequality was shown in
\cite{KR1} in arbitrary metric semigroups -- including Banach spaces as
well as (non-abelian) compact Lie groups. Similarly in \cite{KR2}, the
authors transferred the (sharp) Khinchin--Kahane inequality from Banach
spaces to abelian groups $G$ equipped with a homogeneous length function.
To explore extensions of these results to the non-abelian setting --
e.g., of Lie groups with left-invariant metrics -- we need to first
understand if such objects exist. As explained above, this question was
not answered in the literature; but it is now settled by our main result.

Finally, there may also be a relation to the Ribe program \cite{Naor},
which aims to reformulate aspects of Banach space theory in purely metric
terms. Indeed, from Corollary~\ref{abel} we see that a metric space $X$
is isometric to an additive subgroup of a Banach space if and only if
there is a group structure on $X$ which makes the metric left-invariant
and the length function $\norm{ x } \coloneqq d(1,x)$ homogeneous.
%}}}


\section*{Acknowledgements}

This project is an online collaboration that originated from a blog post
at \url{https://terrytao.wordpress.com/2017/12/16}, following the model
of the ``Polymath'' projects \cite{gn}.  A full list of participants and
their grant acknowledgments may be found at
\begin{center}
\small{\url{http://michaelnielsen.org/polymath1/index.php?title=linear\_norm\_grant\_acknowledgments}}
\end{center}


%{{{1 Section 2 - Key proposition
\section{Key proposition}

The key proposition used to prove Theorem~\ref{class} is the following
estimate, which can treat a somewhat more general class of functions than
homogeneous pseudo-length functions, in which the symmetry hypothesis is
dropped and one allows for an error in the triangle inequality and the
homogeneity property (which is now only claimed for $n=2$).

\begin{proposition}\label{main}
Let $G = (G,\cdot)$ be a group, let $c,c' \in \R$, and let
$\normsymb: G \to \R$ be a function obeying the following axioms:
\begin{itemize}
\item[(i)]  For any $x,y \in G$, one has
\begin{equation}\label{ng}
 \norm{xy} \leq \norm{x}+\norm{y} + c.
\end{equation}
\item[(ii)]  For any $x \in G$, one has
\begin{equation}\label{double}
 \norm{x^2} \geq 2\norm{x} - c'.
\end{equation}
\end{itemize}
Then for any $x,y \in G$, one has
\begin{equation}\label{xyc}
 \norm{ [x,y] } \leq 4 c + 5 c',
\end{equation}
where the commutator $[x,y]$ was defined in \eqref{comdef}.
\end{proposition}

Notably, we neither assume symmetry $\norm{ x^{-1} } = \norm{ x }$, not
even up to a constant, nor $\norm{ e } = 0$; we also allow $\normsymb$ to
take on negative values.  The reader may however wish to restrict
attention to length functions (and set $c=c'=0$) for a first reading of
the arguments below.   The factors of $4,5$ are probably not optimal
here, but the crucial feature of the bound \eqref{xyc} for our
application is that the right-hand side vanishes when $c=c'=0$.

For the remainder of this section, let $G$, $c,c'$, and $\normsymb$
satisfy the hypotheses of the proposition.  Our task is to establish the
bound \eqref{xyc}.  We shall now use \eqref{ng} and \eqref{double}
repeatedly to establish a number of further inequalities relating the
pseudo-lengths $\norm{x}$ of various elements $x$ of $G$, culminating in
\eqref{xyc}.
Many of our inequalities will involve terms that depend on an auxiliary
parameter $n$, but we will be able to eliminate several of them by the
device of passing to the limit $n \to \infty$.  It is because of this
device that we are able to obtain a bound \eqref{xyc} whose right-hand
side is completely uniform in $x$ and $y$.

From \eqref{ng} and induction we have the approximate upper homogeneity
bound
\begin{equation}\label{ng-iter}
 \norm{x^n} \leq n \norm{x} + (n-1) c
\end{equation}
for any natural number $n \geq 1$.  Similarly, from \eqref{double} and
induction one has the lower homogeneity bound
\[
\norm{x^n} \geq n \norm{x} - (n-1) c'
\]
whenever $n$ is a power of two.  It is convenient to rearrange this
latter inequality as
\begin{equation}\label{rearrange}
\norm{x} \leq \frac{\norm{x^n}}{n}  + \frac{n-1}{n} c'.
\end{equation}
This inequality (particularly in the asymptotic limit $n \to \infty$)
will be the principal means by which the hypothesis \eqref{double} is
employed.

We remark that by further use of \eqref{ng-iter} one can also obtain a
similar estimate to \eqref{rearrange} for natural numbers $n$ that are
not powers of two\footnote{For instance, this argument shows that if a
pseudo-length function obeys the doubling condition
$\norm{x^2}=2\norm{x}$ for each $x\in G$, then it is homogeneous; see
also \cite[Lemma 1]{GK}.}, but the powers of two will suffice for the
arguments that follow.

\begin{lemma}[Approximate conjugation invariance]\label{aci}
For any $x,y \in G$, one has
\[
\norm{yxy^{-1} } \leq \norm{ x } + c + c'.
\]
\end{lemma}


\begin{remark}\label{Rinv}
Setting $c=c'=0$, we conclude that any homogeneous pseudo-length function
is conjugation invariant, and thus determines a bi-invariant metric on
$G$. It should not be surprising that this observation is used in the
proof of Theorem~\ref{class}, since it is a simple consequence of that
theorem.
\end{remark}

\begin{proof}[Proof of Lemma~\ref{aci}]
From \eqref{rearrange} (with $x$ replaced by $yxy^{-1}$) one has
\[
\norm{yxy^{-1}} \leq \frac{\norm{y x^n y^{-1}} }{n} + \frac{n-1}{n} c'
\]
whenever $n$ is a power of two.  On the other hand, from \eqref{ng-iter}
and \eqref{ng} one has
\[
\norm{ y x^n y^{-1} } \leq \norm{y} + n \norm{x} + \norm{y^{-1}} + (n+1) c
\]
and thus
\[
\norm{yxy^{-1}} \leq \norm{x} + c + c' + \frac{\norm{y} + \norm{y^{-1}} +
c - c'}{n}.
\]
Sending $n \to \infty$, we obtain the claim.
\end{proof}


\begin{lemma}[Splitting lemma]\label{split}
Let $x,y,z,w \in G$ be such that $x$ is conjugate to both $wy$ and
$zw^{-1}$.  Then one has
\begin{equation}\label{ineq}
\norm{x} \leq \frac{\norm{y} + \norm{ z } }{2}  + \frac{3}{2} (c + c').
\end{equation}
\end{lemma}

\begin{proof}
If we write $x = swys^{-1} = t zw^{-1} t^{-1}$ for some $s,t \in G$, then
from \eqref{rearrange} we have
\begin{align*}
 \norm{x} &\leq \frac{\norm{ x^n x^n } }{2n} + \frac{2n-1}{2n} c' \\
&= \frac{\norm{ s(wy)^n s^{-1}t (zw^{-1})^n t^{-1} } }{2n} +
 \frac{2n-1}{2n} c'
\end{align*}
whenever $n$ is a power of two.  From Lemma~\ref{aci} and \eqref{ng} one
has
\begin{align*}
\norm{(wy)^{k+1} s^{-1}t (zw^{-1})^{k+1} } &= \norm{ w y (wy)^k s^{-1} t
(zw^{-1})^k z w^{-1} } \\
&\leq \norm{ y (wy)^k s^{-1} t (zw^{-1})^k z } + c + c' \\
&\leq \norm{ (wy)^k s^{-1} t (zw^{-1})^k } + \norm{y} + \norm{z} + 3c +
c'
\end{align*}
for any $k \geq 0$, and hence by induction
\[
\norm{(wy)^n s^{-1}t (zw^{-1})^n } \leq \norm{s^{-1} t} + n ( \norm{y} +
\norm{z} + 3c + c' ).
\]
Inserting this into the previous bound for $\norm{x}$ (via two
applications of \eqref{ng}), we conclude that
\[
\norm{x} \leq \frac{\norm{y} + \norm{z}+ 3c + c' }{2} + \frac{\norm{s} +
\norm{s^{-1}t} + \norm{t^{-1}} + 2c}{2n} + \frac{2n-1}{2n} c';
\]
sending $n \to \infty$, we obtain the claim.
\end{proof}


\begin{corollary}
If $x,y \in G$, let $f = f_{x,y}: \Z^2 \to \R$ denote the function
\[
f(m,k) \coloneqq \norm{ x^m [x,y]^k }.
\]
Then for any $m,k \in \Z$, we have
\begin{equation}\label{fmk}
 f(m,k) \leq \frac{f(m-1,k) + f(m+1,k-1)}{2} + 2(c + c').
\end{equation}
\end{corollary}

\begin{proof}
Observe that $x^m [x,y]^k$ is conjugate to both $x (x^{m-1} [x,y]^k)$ and
to $(y^{-1} x^m [x,y]^{k-1} xy) x^{-1}$, hence by \eqref{ineq} one has
\[
\norm{ x^m [x,y]^k } \leq \frac{\norm{ x^{m-1} [x,y]^k } + \norm{ y^{-1}
x^{m} [x,y]^{k-1} xy }}{2} + \frac{3}{2} (c + c').
\]
Since $y^{-1} x^{m} [x,y]^{k-1} xy $ is conjugate to $x^{m+1}
[x,y]^{k-1}$, the claim now follows from Lemma~\ref{aci}.
\end{proof}

We now prove Proposition~\ref{main}. Let $x,y \in G$.  We can
write the inequality \eqref{fmk} in probabilistic form as
\[
f(m,k) \leq {\mathbf E} f\left( \left(m,k-\frac{1}{2}\right) + Y
\left(1,-\frac{1}{2}\right) \right) + 2(c + c')
\]
where $Y = \pm 1$ is a Bernoulli random variable that equals $1$ or $-1$
with equal probability.  The key point here is the drift of $\left(0,
-\frac{1}{2} \right)$ in the right-hand side.
Iterating this inequality, we see that if we let
$n$ be a large power of $2$, then we have
\[
f( 0, n) \leq {\mathbf E} f\left( (Y_1 + \dots + Y_{2n}) \left(1,
-\frac{1}{2}\right) \right) + 4(c + c') n,
\]
where $Y_1,\dots,Y_{2n}$ are independent copies of $Y$ (so in particular
$Y_1+\dots+Y_{2n}$ is an even integer).

From \eqref{ng} and \eqref{ng-iter} one has the inequality
\begin{eqnarray*}
f(m,k) & \leq &
|m| \left(\max( \norm{x}, \norm{x^{-1}})+c\right)\\
& & + |k| \left(\max(
\norm{[x,y]}, \norm{[x,y]^{-1}})
+ c\right) + \norm{e}
\end{eqnarray*}
(the final term being needed to deal with the $m=k=0$ case).  We conclude
that
\[
f\left( (Y_1 + \dots + Y_{2n}) \left(1, -\frac{1}{2}\right) \right) \leq
A |Y_1 + \dots + Y_{2n}| + \norm{e}
\]
where $A$ is a quantity independent of $n$; more explicitly, one can take
\[
A \coloneqq \max\left( \norm{x}, \norm{x^{-1}}\right) + \frac{1}{2}
\max\left(\norm{[x,y]}, \norm{[x,y]^{-1}}\right)  + \frac{3}{2} c.
\]
Taking expectations, since the random variable $Y_1+\dots+Y_{2n}$ has
mean zero and variance $2n$, we see from the Cauchy--Schwarz inequality
(or Jensen's inequality) that
\[
{\mathbf E}|Y_1 + \dots + Y_{2n}|  \leq \left( {\mathbf E}|Y_1 + \dots +
Y_{2n}|^2 \right)^{1/2} = \sqrt{2n}
\]
and hence
\[
f(0,n) \leq A \sqrt{2n} + \norm{e} + 4 (c + c') n.
\]
But from \eqref{rearrange} we have
\[
\norm{[x,y]} \leq \frac{f(0,n)}{n}  + \frac{n-1}{n} c'.
\]
Combining these two bounds and sending $n \to \infty$, we obtain
Proposition~\ref{main}.


\begin{remark}
As discussed in the footnote preceding Lemma \ref{aci}, if $c=0$ then the
condition $c'=0$, or equivalently $\norm{x^2} = 2 \norm{x}$ for all $x \in
G$, is in turn equivalent to $\norm{x^k} = k \norm{x}$ for all $x \in G$,
for any $k \geq 2$. One can thus ask if \eqref{double} can be replaced by
a similar condition of the form
\[
\norm{x^k} \geq k \norm{x} - c'
\]
for some $k \geq 2$. We claim that such a condition essentially implies
\eqref{double}. Indeed, using the triangle inequality for $x^k =
(x^2)^{\lfloor k/2 \rfloor} x^\delta$ where $\delta=0,1$, we obtain:
\[
k \norm{x} - c' \leq \norm{x^k} \leq O(k) \norm{x^2} + c \, O(k),
\]
which in turn implies $\norm{x^2} \geq O(1) \norm{x} - O(1) c - c'/k$.
\end{remark}


\begin{remark}
One can deduce a `local' version of Proposition \ref{main} as follows:
notice that the constants $c,c'$ can be described in terms of $\normsymb$
from \eqref{ng} and \eqref{double}, to yield
\begin{equation}\label{9bound}
\norm{ [x,y] } \leq 4 \left( \sup_{z,w \in G} \norm{zw} - \norm{z} -
\norm{w} \right) + 5 \left( \sup_{z \in G} 2 \norm{z} - \norm{z^2}
\right),
\end{equation}
for any group $G$ and function $\normsymb : G \to \R$, and any $x,y \in
G$.
(Both sides are zero when $G$ is a Banach space and $\normsymb$ is the
norm, so equality is obtained in that case.)
Now it is enough to consider $G = \langle x,y\rangle$ without loss of
generality, which may lead to a better bound on $\norm{ [x,y] }$ than
taking the suprema over all of $G$.
Notice also that at least one of the constants $c,c'$ must be
non-negative, from \eqref{ng} with $x=y$ and \eqref{double}:
\begin{equation}\label{Enonneg}
\max \{ c, c' \} \geq \max\left\{ \norm{e^2} - 2 \norm{e}, 2 \norm{e} -
\norm{e^2}\right\} = \max \{ -\norm{e}, \norm{e} \} \geq 0.
\end{equation}
In fact, this reasoning and our results imply that the only way to get
$c=c'=0$ on the right-hand side of \eqref{9bound} is when $\normsymb$
arises from pulling back the norm of a Banach space $\mathbb{B}$ along a
group homomorphism $G\to\mathbb{B}$, or equivalently along a group
homomorphism from the torsion-free abelianization of $G$ to $\mathbb{B}$.
\end{remark}
%}}}


%{{{1 Section 3 - Proof of main theorem
\section{Proof of main theorem}

With Proposition~\ref{main} in hand, it is not difficult to conclude the
proof of Theorem~\ref{class}.  Suppose that $G$ is a group with a
homogeneous pseudo-length function $\normsymb: G \to [0,+\infty)$.
Applying Proposition~\ref{main} with $c=c'=0$, we conclude that
$\norm{[x,y]}=0$ for all $x,y \in G$, thus by the triangle inequality
$\normsymb$ vanishes on the commutator subgroup $[G,G]$, and therefore
factors through the abelianization $G_{\mathrm{ab}} \coloneqq G/[G,G]$ of
$G$.  (This already establishes part of one implication of
Corollary~\ref{abel}.)  Factoring out by $[G,G]$ like this, we may now
assume without loss of generality that $G$ is abelian.  To
reflect this, we now use additive notation for $G$, thus for instance
$\norm{nx} = |n| \norm{x}$ for each $x \in G$ and $n \in \Z$, and one can
also view $G$ as a module over the integers $\Z$.

At this point we repeat the arguments in \cite[Theorem B]{KR2}, which
treated the case when $G$ was separable, though it turns out that this
separability hypothesis is unnecessary.

If $x$ is a torsion element of $G$ (thus $nx=0$ for some $n$), then the
homogeneity condition forces $\norm{x}=0$.  Thus $\normsymb$ vanishes on
the torsion subgroup of $G$; factoring out by this subgroup, we may thus
assume without loss of generality that $G$ is not only abelian, but is
also torsion-free.

We can view $G$ as a subgroup of the $\Q$-vector space $G \otimes_\Z \Q$,
the elements of which can be formally expressed as $\frac{1}{n} x$ for
natural numbers $n$ and elements $x \in G$ (with two such expressions
$\frac{1}{n} x, \frac{1}{m} y$ identified if and only if $mx = ny$, and
the $\Q$-vector space operations defined in the obvious fashion); the
fact that this is well defined as a $\Q$-vector space follows from the
hypotheses that $G$ is abelian and torsion-free.  We can then define a
seminorm $\| \, \|_\Q: G \otimes_\Z \Q \to [0,+\infty)$ by setting
\[
\left\|\frac{1}{n} x\right\|_\Q \coloneqq \frac{1}{n} \norm{x}
\]
for any $x \in G$ and natural number $n$; the linear growth condition
ensures that $\|\,\|_\Q$ is well-defined.  It is not difficult to verify
the homogeneity condition
\[
\| q x \|_\Q = |q| \|x\|_\Q
\]
and the triangle inequality
\[
\|x+y\|_\Q \leq \|x\|_\Q + \|y\|_\Q
\]
for all $x,y \in G \otimes_\Z \Q$ and $q \in \Q$, so that $\| \, \|_\Q$ is
indeed a seminorm over the $\Q$-vector space $G \otimes_\Z \Q$.

The $\Q$-vector space $G \otimes_\Z \Q$ may be embedded in turn into the
$\R$-vector space $G \otimes_\Z \R = (G \otimes_\Z \Q) \otimes_\Q \R$,
elements of which take the form $\alpha_1 x_1 + \dots + \alpha_k x_k$ for
some $x_1,\dots,x_k \in G \otimes_\Z \Q$ and $\alpha_1,\dots,\alpha_k \in
\R$.  We identify two such elements $\alpha_1 x_1 + \dots + \alpha_k
x_k$, $\beta_1 y_1 + \dots + \beta_l y_l$ if and only if there are
$\Q$-linearly independent elements $z_1,\dots,z_m \in G \otimes_\Z \Q$
and coefficients $a_{in}, b_{jn} \in \Q$ for $i=1,\dots,k$,
$j=1,\dots,l$, $n=1,\dots,m$ such that $x_i = \sum_{n=1}^m a_{in} z_n$
for each $i=1,\dots,k$, $y_j = \sum_{n=1}^m b_{jn} z_n$ for each
$j=1,\dots,l$, and $\sum_{i=1}^k \alpha_i a_{in} = \sum_{j=1}^l \beta_j
b_{jn}$ for each $n=1,\dots,m$; one can check that this criterion does not
depend on the precise choice of linearly independent elements
$z_1,\dots,z_m$, so long as their span contains all of the
$x_1,\dots,x_k, y_1,\dots,y_l$.\footnote{Stated slightly differently, let
$z_1, \dots, z_m$ denote a $\Q$-basis of the $\Q$-span of the $x_i, y_j$,
and define the vector $\bm{z} \coloneqq (z_1, \dots, z_m)^T$ by abuse of
notation; similarly define $\bm{x}, \bm{y}, \bm{\alpha}, \bm{\beta}$.
Then $\bm{x} = A \bm{z}, \ \bm{y} = B \bm{z}$ for matrices
$A \in \Q^{k \times m}, B \in \Q^{l \times m}$, and
we identify $\bm{\alpha}^T \bm{x} = \bm{\beta}^T \bm{y}$ if and only if
$\bm{\alpha}^T A = \bm{\beta}^T B$.}
It is easy to verify\footnote{One subtle point is that if $G$ is already
embedded inside an existing $\R$-vector space $V$, then the vector space
$G \otimes_\Z \R$ constructed here is not necessarily a subspace of that
vector space (although there will be a canonical linear map from $G
\otimes_\Z \R$ to $V$).  For instance, if $G = \langle 1, \sqrt{2}
\rangle \subset \R$, then while $G$ is contained in a one-dimensional
$\R$-vector space, $G \otimes_\Z \R$ is in fact two-dimensional (with the
generators $1$ and $\sqrt{2}$ of $G$ forming a $\R$-basis of this space);
the point is that the vector space operations on $G \otimes_\Z \R$ do not
agree in this case with those of the ambient vector space $\R$, for
instance in $G \otimes_\Z \R$ the second generator $\sqrt{2}$ is not
equal to the scalar multiple of the first generator $1$ by $\sqrt{2}$.
To avoid confusion, we recommend to the reader that they view $G$ as an
abstract group that is not already embedded into an existing $\R$-vector
space.} that $G \otimes_\Z \R$ has the structure of an $\R$-vector space
that contains $G \otimes_\Z \Q$ (and hence $G$).  We can extend the
seminorm $\|\,\|_\Q$ to $G \otimes_\Z \R$ by setting
\[
\|\alpha_1 x_1 + \dots + \alpha_k x_k\|_\R \coloneqq
\lim_{(q_1,\dots,q_k) \to (\alpha_1,\dots,\alpha_k)} \| q_1 x_1 + \dots +
q_k x_k\|_\Q
\]
where $(q_1,\dots,q_k)$ is restricted to $\Q^k$.  For any given
$\alpha_1,\dots,\alpha_k,x_1,\dots,x_k$ it is easy to see using the
triangle inequality and homogeneity that the limit exists.  We still need
to verify that the seminorm is well-defined, in the sense that
\[
\|\alpha_1 x_1 + \dots + \alpha_k x_k\|_\R  = \| \beta_1 y_1 + \dots +
\beta_l y_l \|_\R
\]
whenever $\alpha_1 x_1 + \dots + \alpha_k x_k = \beta_1 y_1 + \dots +
\beta_l y_l$.  By the triangle inequality it suffices to verify this at
the origin, or more precisely that
\[
\lim_{(q_1,\dots,q_k) \to (\alpha_1,\dots,\alpha_k)} \| q_1 x_1 + \dots +
q_k x_k\|_\Q = 0
\]
whenever $x_1,\dots,x_k \in G \otimes_\Z \Q$ and $\alpha_1,\dots,\alpha_k
\in \R$ are such that $\alpha_1 x_1 + \dots + \alpha_k x_k = 0$.  To
prove this one can induct on $k$, the case $k=0$ being trivial.  If one
of the $x_1,\dots,x_k$ is a $\Q$-linear combination of the others, then
we can eliminate it and invoke the induction hypothesis for $k-1$, so we
may assume without loss of generality that the $x_1,\dots,x_k$ are
linearly independent over $\Q$.  But then, by construction of $G
\otimes_\Z \R$, the condition $\alpha_1 x_1 + \dots + \alpha_k x_k = 0$
forces $\alpha_1=\dots=\alpha_k=0$, and the claim now follows from the
triangle inequality and homogeneity.

By construction, $\| \, \|_\R: G \otimes_\Z \R \to
[0,+\infty)$ obeys the homogeneity condition
\[
\| \alpha x \|_\R = |\alpha| \|x\|_\R
\]
and the triangle inequality
\[
\|x+y\|_\R \leq \|x\|_\R + \|y\|_\R
\]
for all $x,y \in G \otimes_\Z \R$ and $\alpha \in \R$; thus $\|\,\|_\R$
is a seminorm on $G \otimes_\Z \R$.   By the usual metric completion
device of taking formal Cauchy sequences and factoring out the elements
that converge to zero, we can find a real Banach space $\mathbb{B} =
(\mathbb{B},\|\,\|_{\mathbb{B}})$ and an $\R$-linear map $\phi: G
\otimes_\Z \R \to \mathbb{B}$ such that $\|x\|_\R = \| \phi(x)
\|_{\mathbb{B}}$ for each $x \in G \otimes_\Z \R$.
(In fact, as observed in \cite{KR2}, one can directly take the metric
space completion of $G \otimes_\Z \Q$ without first passing to $G
\otimes_\Z \R$.)
It is now clear that $\phi$ restricts to a homomorphism from $G$ to
$\mathbb{B}$ such that $\norm{x} = \| \phi(x) \|_{\mathbb{B}}$ for each
$x \in G$.  Finally, if $\normsymb$ is a length function, then $\|
\phi(x) \|_{\mathbb{B}}$ must be non-zero for each $x \neq e$, which
implies that the kernel of $\phi$ is trivial and hence $\phi$ is
injective.

This concludes the proof of Theorem~\ref{class}.  Since the homomorphism
$\phi: G \to \mathbb{B}$ can only be injective for abelian torsion-free
$G$, we obtain the ``only if'' portion of Corollary~\ref{abel}.
Conversely, if a group $G$ is abelian and torsion-free, by the above
constructions it embeds into a real vector space $\mathbb{B} \coloneqq G
\otimes_\Z \R$; now by Zorn's lemma $\mathbb{B}$ has a norm (e.g.,
consider the $\ell^1$ norm with respect to a Hamel basis of
$\mathbb{B}$), which restricts to the desired homogeneous length function
on $G$. We remark, as also mentioned in \cite{KR2}, that $G \otimes_\Z
\R$ is the construction of the smallest, `enveloping' vector space
containing a copy of the abelian, torsion-free group $G$.

\begin{remark}
The above arguments also show that homogeneous pseudo-length functions on
$G$ are in bijection with seminorms on the real vector space
$G_{\mathrm{ab},0} \otimes_\Z \R$, where $G_{\mathrm{ab},0}$ denotes the
torsion-free abelianization of $G$.
\end{remark}
%}}}


\section{Further remarks and results}\label{remarks-sec}

One could ask if the homogeneity condition can be relaxed to an
asymptotic homogeneity condition
\[
\lim_{n \to \infty} \frac{\norm{ x^n }}{ n } = \norm{ x }, \ \text{ for
each }\ x \in G.
\]
In fact this condition is equivalent to homogeneity,
since for all positive integers $m \in \N$, the sequence
$\{ \norm{ x^{mn} } / (mn) : n \geq 1 \}$ would converge
as $n \to \infty$ to both $\norm{ x }$ and $\norm{ x^m } / m$.

However, if we weaken any of several conditions in the above
formulations, then examples of such structures do, in fact, exist. Here
we mention several such cases and discuss other related problems.


%{{{1 Section 4.1 - Monoids and embeddings
\subsection{Monoids and embeddings}

Our first weakening is to replace `groups' by the more primitive
structures `monoids' or `semigroups'. In this case, Robert Young (private
communication) described to us non-abelian monoids with homogeneous,
bi-invariant length functions:
consider the free monoid $\mathrm{FMon}(X)$ on any alphabet $X$ of size
at least $2$, with the `edit' distance $d(v,w)$ between strings $v,w \in
\mathrm{FMon}(X)$ being the least number of single generator insertions
and deletions to get from $v$ to $w$. The triangle inequality and
positivity are easily verified, while homogeneity of the corresponding
length function $\norm{x} \coloneqq d(e,x)$ is trivial. Moreover, the
metric $d(\cdot,\cdot)$ turns out to be bi-invariant:
\[
d(gxh,gyh) = d(x,y) \ \text{ for all }\ g,h,x,y \in \mathrm{FMon}(X).
\]
This specializes to left- and right-invariance upon taking $g \in X$ and
$h = e$, or  $h \in X$ and $g = e$, respectively.

Note moreover that $\mathrm{FMon}(X)$ embeds into $\mathrm{FGp}(X)$ (the
free group generated by $X$ and $X^{-1}$, the collection of symbols
defined to be inverses of elements of $X$), and in particular, is
cancellative. While this trivially addresses the embeddability issue,
notice that a more refined version of embeddability fails. Namely, by our
main theorem, $\mathrm{FMon}(X)$ does not embed into any group in the
category $\mathcal{C}_{\rm{bi-inv,\ hom}}$ with cancellative semigroups
with homogeneous bi-invariant metrics as objects and isometric semigroup
maps as morphisms. Thus, one may reasonably ask what is a sufficiently
small category in which the embeddability works.

\begin{proposition}\label{Pisometry}
Let $\mathcal{C}_{\rm{bi-inv}}$ denote the category whose objects are
cancellative semigroups with bi-invariant metrics, and morphisms are
isometric semigroup maps. Then $\mathrm{FMon}(X)$ embeds isometrically
into $\mathrm{FGp}(X)$ in $\mathcal{C}_{\rm{bi-inv}}$.
\end{proposition}

\begin{proof}
From above, $\mathrm{FMon}(X)$ is an object of
$\mathcal{C}_{\rm{bi-inv}}$; denote the metric by $d_{FM}$. One can check
that $d_{FM}(w,w')$ equals the difference $l(w) + l(w')$ (the sum of the
lengths in $X$) minus twice the length of the longest common (possibly
non-contiguous) substring in $w,w'$.

We next claim $\mathrm{FGp}(X)$ is also an object of
$\mathcal{C}_{\rm{bi-inv}}$. Namely, for a word $w=x_1x_2\cdots x_m$ in
the free group, we consider \emph{non-crossing matchings} in $w$, i.e.,
sets $M$ of pairs of letters in $\{1,2,\dots m\}$ such that the following
hold.
\begin{itemize}
\item If $(i, j)\in M$, then $i< j$ and  $x_j=x_i^{-1}$.
\item If $(i, j)\in M$ and $(k, l)\in M$, then either $(i, j) = (k, l)$
or $i$, $j$, $k$ and $l$ are distinct.
\item If $i<k<j<l$ and $(i, j)\in M$, then $(k, l)\notin M$.
\end{itemize}

Given a matching $M$ as above, consider the set $U=U(M)$ of indices $k$,
$1\leq k\leq m$ which are not part of a pair in $M$. Define the
\emph{deficiency} of the matching $M$ as the cardinality of the set
$U(M)$, and define the length $\norm{w}$ of the word $w$ as the infimum
of the deficiency over all non-crossing matchings in $w$.
This length was previously studied in \cite{gadgil}, including checking
that it is well-defined on all of $\mathrm{FGp}(X)$; moreover, $\norm{w}$
equals the smallest number of conjugates of elements in $X \sqcup X^{-1}$
whose product is $w$. Now define $d_{FG}(w,w') \coloneqq \norm{w^{-1}
w'}$. It is easy to see that $\normsymb$ is a conjugacy invariant length
function.

We claim that $d_{FG} \equiv d_{FM}$ on $\mathrm{FMon}(X)$, which proves
the result. It is easy to show that if two words in $\mathrm{FMon}(X)$
differ by a single insertion or deletion, then their distance in
$\mathrm{FGp}(X)$ is at most one, hence exactly one. In the other
direction, we claim that a non-crossing matching on $w^{-1}w'$, with $w$
and $w'$ containing only positive generators (in $X$), is just a
`rainbow', i.e.~nested arches with one end in $w^{-1}$ and the other in
$w'$. But then $d_{FG}(w,w')$ equals $l(w) + l(w')$ minus twice the
length of a common substring, which is maximal by the minimality of the
deficiency. Hence $d_{FG}(w,w') = d_{FM}(w,w')$, completing the proof.
\end{proof}
%}}}


%{{{1 Section 4.2 - Quasi-pseudo-length functions via quasimorphisms
\subsection{Quasi-pseudo-length functions via quasimorphisms}

We see that Proposition~\ref{main} is also of interest for $c > 0$, as
follows. A \emph{quasi-morphism} on a group $G$ is a map $f : G \to \R$
whose \emph{defect} is bounded,
\[
	D(f) \coloneqq \sup_{x,y\in G} | f(xy) - f(x) - f(y) | < +\infty.
\]
Let us use the term \emph{quasi-pseudo-length function} for a function
$\normsymb : G \to [0,+\infty)$ which satisfies the relaxed triangle
inequality~\eqref{ng} for some $c$. Then every quasi-morphism as above
induces a quasi-pseudo-length function by setting
\begin{equation}\label{quasimortonorm}
	\norm{ x } \coloneqq | f(x) |,
\end{equation}
where we can take $c = D(f)$.

A quasi-morphism is \emph{homogeneous} if $f( x^n ) = n f( x )$ for $n
\in \N$. Replacing any quasi-morphism $f$ by
$\lim_{n \to \infty} f( x^n) / n$ results in a homogeneous
quasi-morphism. (Note that this limit exists by the de Bruijn--Erd\"os
generalization \cite[Theorem 23]{BE} of Fekete's subadditive lemma.)
A homogeneous quasi-morphism defines a quasi-pseudo-length function that
is homogeneous.\footnote{Of course, homogeneous quasi-morphisms must
satisfy
\[
\sup_{n \geq 1} |f(e) - f(g^n) - f(g^{-n})| = \sup_{n \geq 1} |f(e) -
n(f(g) + f(g^{-1}))| < \infty,
\]
whence $f(g^{-1}) = -f(g)$ for all $g \in G$. On the other hand, if $f
\not\equiv 0$ then the corresponding homogeneous quasi-pseudo-length
function $\normsymb \coloneqq |f|$ never has bounded defect, since if
$f(g) \neq 0$ then $\inf_{n \geq 0} \norm{e} - \norm{g^n} - \norm{g^{-n}}
= \inf_{n \geq 0} -2n \norm{g} = -\infty$. So in a sense, the two
concepts are `morally different'.}
In this case, Proposition~\ref{main} makes a rather trivial statement: a
homogeneous quasi-morphism is bounded on commutators,
\[
	| f([x,y]) | \leq 4 D(f).
\]
In fact, as observed in \cite[Lemme~1.1]{Br}, one can improve the
constant $4$ to $3$ in this case.

Nevertheless, quasi-morphisms can be utilized to construct interesting
quasi-pseudo-length functions, for example satisfying homogeneity on
specific commutators. The following quasi-morphism is due to
Brooks~\cite[Section~2]{Fu}.
For a given word $w$ in the free group $\free{2}$, written in reduced
form, let $f_w : \free{2} \to \R$ be the function which assigns to every
other $g\in \free{2}$, also written in reduced form, the maximum number
of times such that $w$ occurs in $g$ without overlaps, minus the
analogous number of times that $w^{-1}$ can maximally occur in $g$. Since
$f_w(w^n) = n f_w(w)$, homogenizing this quasi-morphism and taking the
absolute value results in a quasi-pseudo-length function that is
homogeneous on the powers of $w$. For example with $w$ being the
commutator of the generators of $\free{2}$, we see that although the
quasi-pseudo-length function must be bounded on commutators by
Proposition~\ref{main}, it can nevertheless grow linearly on the powers
of a fixed commutator.
%}}}


%{{{1 Section 4.3 - Finite balls in free groups
\subsection{Finite balls in free groups}

From Proposition~\ref{main} and a standard compactness argument, we can
establish the following local version of the theorem.

\begin{theorem}
For any $\eps>0$ there exists $R \geq 4$ with the following property: if
$a,b$ are two elements of a group $G$, $B_{a,b}(R) \subset G$ is the
collection of all words in $a,b,a^{-1},b^{-1}$ of length at most $R$
\textup{$($}so in particular $B_{a,b}(R)$ contains $[a,b]$\textup{$)$},
and the map $\normsymb: B_{a,b}(R) \to [0,+\infty)$ is a ``local
pseudo-length function'' which obeys the triangle inequality
\begin{equation}\label{tri}
 \norm{ xy } \leq \norm{x} + \norm{y}
\end{equation}
whenever $x,y,xy \in B_{a,b}(R)$, with equality when $x=y$, then one has
\[
\norm{ [a,b] } \leq \eps ( \norm{a} + \norm{b} ).
\]
\end{theorem}

\begin{proof}
By pulling back to the formal group generated by $a$ and $b$, we may
assume without loss of generality that $G$ is the free group generated by
$a$ and $b$.  Without loss of generality we may also normalize
$\norm{a}+\norm{b}=1$.  If the claim failed, then one could find a
sequence $R_n \to \infty$ and local pseudo-length functions $\normsymb_n:
B_{a,b}(R_n) \to [0,+\infty)$ such that $\normn{a}+\normn{b} = 1$, but
that $\normn{[a,b]} \geq \eps$.  By the Arzela--Ascoli theorem, we can
pass to a subsequence that converges pointwise to a homogeneous
pseudo-length function $\normsymb: G \to [0,+\infty)$ such that
$\norm{[a,b]} \geq \eps$, which contradicts Proposition~\ref{main}.
\end{proof}

\begin{remark}
By carefully going through the arguments in the previous section,
choosing $n$ to be various small powers of $R$ instead of sending $n$ to
infinity, one can extract an explicit value of $R$ of the form $R = C
\eps^{-A}$ for some absolute constants $C, A > 0$; we leave the details
to the interested reader.
\end{remark}


On the other hand, for any finite $R$ one can construct local
pseudo-length functions $\normsymb: B(0,R) \to [0,+\infty)$ such that
$\norm{x} > 0$ for all $x \in B(0,R) \setminus \{e\}$.  One construction
is as follows. Any two matrices $U_a, U_b\in SO(3)$ define a
representation $x \mapsto U_x$ of the free group $\free{2}$ in the
obvious fashion.  Every $U_x$ is then a rotation around some axis in
$\R^3$ by some angle $0 \leq \theta \leq \pi$ in one of the two
directions around that axis; if $U_a$ and $U_b$ are sufficiently close to
the identity, then the angle of rotation of $U_x$ will be at most $\pi/2$
for all $x \in B(0,R)$. We set $\norm{x}$ to be equal to this angle of
rotation.  Also, if $U_a,U_b$ are chosen generically, the representation
will be faithful (this follows from the dominance of word maps on simple
Lie groups such as $SO(3)$, see \cite{borel}), so $\norm{x}>0$ for any
non-identity $x$.  From the triangle inequality for angles we thus have
\eqref{tri} whenever $x,y,xy \in B(0,R)$, with equality when $x=y$.  Note
that as one sends $R \to \infty$, the local pseudo-length functions
constructed here converge to zero pointwise, so in the limit we do not
get any counterexample to the main theorem.
%}}}


%{{{1 Section 4.4 - Quasi-norms and commutator lengths
\subsection{Quasi-norms and commutator lengths}

Our main result is  a classification theorem for homogeneous
pseudo-length functions. We will see next that even in the case of
homogeneous quasi-length functions, i.e., where we replace $\norm{xy}
\leq \norm{x} + \norm{y}$ by $\norm{xy} \leq \norm{x} + \norm{y} + c$ for
a fixed constant $c\geq 0$, we get non-trivial restrictions on these
functions which give the main result in the case when $c=0$. Further, we
see that the restrictions we obtain are essentially the strongest
possible.

Recall that the \emph{commutator length} $\mathrm{cl}(g)$ of a word in
$[G, G]$ is the length $k$ of the shortest expression $g = [a_1,
b_1]\cdot[a_2, b_2]\cdots[a_k, b_k]$ of $g$ as a product of commutators.
The \emph{stable commutator length} is defined as $\lim_{n\to\infty}
\mathrm{cl}(g^n)/n$, where the limit exists by sub-additivity of the
function $n\mapsto \mathrm{cl}(g^n)$.

Then Proposition~\ref{main}, together with \eqref{ng}, \eqref{rearrange},
easily imply the following estimates:

\begin{proposition}
Let $\normsymb$, $c$, $c'$ be as in Proposition~\ref{main}. Then for
$g\in [G, G] \setminus \{e\}$,
$\norm{g}\leq 5(c+c')\, \mathrm{cl}(g) - c$ and $\norm{g}\leq 5(c+c')\,
\mathrm{scl}(g) + c'$.
\end{proposition}


Observe that in various cases where we do not have a (quasi-)homogeneous
quasi-norm, we get much weaker bounds. For example, consider the word
$[a^k, b^m]$ in the free group $\free{2}$, generated by $a$ and $b$, for
some integers $k$ and $m$.
\begin{itemize}
\item The norm of such an element with respect  to the word metric is
$2(|k| + |m|)$.

\item If we have a length function $\normsymb$ which is symmetric and
conjugation-invariant (but not assumed to be homogeneous), then we have
the bound $\norm{[a^k, b^m]} \leq 2 \min (|k|\, \norm{a}, |m|\,
\norm{b})$.  Furthermore, there are conjugation-invariant length
functions for which these inequalities hold with equality.

Namely, fix positive real numbers $\norm{a}$ and $\norm{b}$; now given a
word $w = x_1 \cdots x_m$, consider the set $U=U(M)$ of indices $k$,
$1\leq k\leq m$ which are not part of a pair in a matching $M$; here $M$
is defined as in the proof of Proposition \ref{Pisometry}. Now define the
deficiency of $M$ as the sum $\sum_{i\in U}\norm{x_i}$, and as in
\textit{loc.~cit.}, define $\norm{w}$ to be the infimum of the deficiency
over all non-crossing matchings in $w$. Then $\normsymb$ is a
conjugation-invariant length function.
Further, $\norm{[a^k, b^m]} \geq 2 \min (|k|\, \norm{a}, |m|\, \norm{b})$
as, for any matching $M$ for $w=[a^k, b^m]$, if some pair $(i, j)$
corresponds to letters $a$ and $a^{-1}$, then no pair corresponds to
letters $b$ and $b^{-1}$ and conversely. Further, it is easy to find a
matching for which the deficiency is $\min (|k|\, \norm{a}, |m|\,
\norm{b})$.  On the other hand, $\normsymb$ is not homogeneous; for
instance, $\norm{[a,b]} = 2$ and $\norm{[a,b]^3} = 4$.  Similarly, we
have $\norm{[a^k,b^k]}=2|k|$ and $\norm{[a^k,b^k]^3} \leq 4|k|$, which
demonstrates that $2\norm{x} - \norm{x^2}$ is unbounded (as must be the
case, according to \eqref{9bound}).

\item On the other hand, the function $\normsymb$ associating to each
word the length of its cyclically reduced form is homogeneous but not a
quasi-length function. For this we have $\norm{[a^k, b^m]} = 2(|k| +
|m|)$.
\end{itemize}

Observe that all of the bounds on $\norm{[a^k, b^m]}$ here become
unbounded as $k,m \to \infty$.  This should be compared with Proposition
\ref{main}, which establishes a bound $\norm{[a^k,b^m]} \leq 4c+5c'$ that
is uniform in $k$ and $m$ for any function $\normsymb$ satisfying the
hypotheses of that proposition.

Now we show that the bound in Proposition \ref{main} is essentially sharp
for every finitely generated group $G$.
By assumption, also the abelianization $G_{\mathrm{ab}} \coloneqq
G/[G,G]$ is finitely generated. For simplicity of notation, assume
$G_{\mathrm{ab}}$ is torsion-free. Let $x_1$, $x_2$, \dots, $x_k$ be the
generators of $G_{\mathrm{ab}} \cong \Z^k$ and let $\varphi: G \to
G_{\mathrm{ab}}$ be the abelianization map. Let $\alpha_i : G \to \Z$ for
$i = 1,\ldots, k$ be the homomorphisms such that for all $g\in G$,
\[
\varphi(g) = \sum_{i=1}^k \alpha_i(g)x_i.
\]
Choose and fix elements $y_i\in G$ such that $\varphi(y_i) = x_i$. We
define a function $\psi: G \to \mathbb{R}$ by
\[
\psi(g) = \mathrm{cl}(g y_1^{-\alpha_1(g)}y_2^{-\alpha_2(g)}\cdots
y_k^{-\alpha_k(g)}).
\]

Repeatedly using relations of the form $\mathrm{cl}(ba) =
\mathrm{cl}(ba[a^{-1}, b^{-1}]) \leq \mathrm{cl}(ab) + 1$, we see that
\[
\psi(g_1g_2)\leq \psi(g_1) + \psi(g_2) + (2k + 1),
\]
i.e., $\psi$ is a quasi-pseudo-length function. Further, by similar
arguments we see that
\[
\vert\psi((xy)^n) - \psi(x^ny^n)\vert \leq n,
\]
from which we deduce that on replacing $\psi$ by its homogenization
(via \cite[Theorem 23]{BE}), we still get a quasi-pseudo-length function.

Therefore, we can define functions satisfying Proposition~\ref{main}
which are homogeneous by considering a length function
$\normsymb_{G_{\mathrm{ab}} }$ on $G_{\mathrm{ab}}$ and taking the linear
combination
\[
\norm{g} \coloneqq \normsymb_{G_{\mathrm{ab}} }(\varphi(g)) +
\frac{c}{2k+1} \psi(g).
\]
This gives a homogeneous function satisfying the hypotheses of
Proposition~\ref{main} while being strictly positive on all nontrivial
elements.
%}}}


%{{{1 Bibliography
\begin{thebibliography}{10}

\bibitem{Br}
C.~Bavard, {\em Longueur stable des commutateurs}, L'Enseignement
  Math\'ematique, 37(1--2):109--150, 1991.

\bibitem{borel}
A.~Borel, \emph{On free subgroups of semisimple groups}, L'Enseignement
  Math\'ematique, 29(1--2):151--164, 1983.

\bibitem{BE}
N.G.~de Bruijn and P.~Erd\"os, {\em Some linear and some quadratic
  recursion formulas. II},
  \href{http://dx.doi.org/10.1016/S1385-7258(52)50021-0}{Indagationes
  Mathematicae}, 55:152--163, 1952.

\bibitem{CSC}
F.~Cabello S\'anchez and J.M.F.~Castillo, {\em Banach space techniques
  underpinning a theory for nearly additive mappings},
  \href{https://www.impan.pl/en/publishing-house/journals-and-series/dissertationes-mathematicae/all/404}{Dissertationes
  Mathematicae} (Rozprawy Matematyczne), 404:1--73, 2002.

\bibitem{Fu}
K.~Fujiwara, {\em Quasi-homomorphisms on mapping class groups},
Handbook of Teichm\"uller theory. Vol. II, 241--269,
IRMA Lect. Math. Theor. Phys., 13, Eur. Math. Soc., Z\"urich, 2009.

\bibitem{gadgil}
S.~Gadgil, {\em Watson--Crick pairing, the Heisenberg group and Milnor
  invariants}, \href{http://dx.doi.org/10.1007/s00285-008-0223-x}{Journal
  of Mathematical Biology} 59(1), Article:123, 2009.
	
\bibitem{GK}
Z.~Gajda and Z.~Kominek, {\em On separation theorems for subadditive and
  superadditive functionals},
  \href{http://pldml.icm.edu.pl/pldml/element/bwmeta1.element.bwnjournal-article-smv100i1p25bwm}{Studia
  Mathematica}, 100(1):25--38, 1991.

\bibitem{gn}
T.~Gowers and M.~Nielsen, {\em Massively collaborative mathematics},
  \href{https://www.nature.com/articles/461879a}{Nature},
  461(7266):879--881, 2009.

\bibitem{KR1}
A.~Khare and B.~Rajaratnam, {\em The Hoffmann-J{\o}rgensen inequality in
  metric semigroups}, \href{http://dx.doi.org/10.1214/16-AOP1160}{Annals
  of Probability}, 45(6A):4101--4111, 2017.

\bibitem{KR2}
A.~Khare and B.~Rajaratnam, {\em The Khinchin--Kahane inequality and
  Banach space embeddings for metric groups}, preprint
  (\href{http://arxiv.org/abs/1610.03037}{arXiv:1610.03037}), 2016.

\bibitem{Koba}
S.~Kobayashi, {\em Hyperbolic Manifolds and Holomorphic Mappings},
World Scientific, 1970.

\bibitem{LT}
M.~Ledoux and M.~Talagrand, {\em Probability in Banach Spaces
  (Isoperimetry and Processes)}, Ergebnisse der Mathematik und ihrer
  Grenzgebiete, Springer-Verlag, Berlin-New York, 1991.

\bibitem{Milnor}
J.~Milnor, {\em Curvatures of left invariant metrics on lie groups},
  \href{http://dx.doi.org/10.1016/S0001-8708(76)80002-3}{Advances in
  Mathematics}, 21(3):293--329, 1976.

\bibitem{Naor}
A.~Naor, {\em An introduction to the Ribe program},
\href{http://dx.doi.org/10.1007/s11537-012-1222-7}{Japan Journal of
Mathematics}, 7:167--233, 2012.

\end{thebibliography}
%}}}

\end{document}
